% Level 2, Week 4 -- While Loops and Game Loops (student workbook)
\chapterweek{4}{While Loops and Game Loops}{A while loop repeats until something changes, and every game is a loop.}

\section{What you will learn}
\begin{itemize}[leftmargin=1.2em]
  \item How a \textbf{while loop} keeps going until something changes.
  \item How to spot and stop an \textbf{infinite loop}.
  \item How to keep asking until the player types a good answer (\textbf{input validation}).
  \item How a \textbf{loop inside a loop} turns a list of strings into pixel art.
\end{itemize}

\section{The big idea: repeat until something changes}
A \texttt{for} loop knows how many times it will go around. A \texttt{while} loop does not: it checks a condition at the top,
and goes around again \emph{while} the condition is \texttt{True}.

\begin{center}
%% alt: A while loop. Start with n = 3. Check: is n bigger than 0? If yes (True), print n, subtract 1 from n and go back to the check. If no (False), leave the loop and print Go!
\begin{tikzpicture}[font=\small,box/.style={draw=black,thick,rounded corners=3pt,align=center,inner sep=6pt},>={Stealth[length=3mm]}]
  \node[box,fill=kcsky!20] (start) at (0,0) {\texttt{n = 3}};
  \node[box,fill=kcyellow!35,diamond,aspect=2.2,inner sep=2pt] (check) at (0,-1.8) {\texttt{n > 0} ?};
  \node[box,fill=kcgreen!20] (body) at (5,-1.8) {\texttt{print(n)}\\\texttt{n = n - 1}};
  \node[box,fill=kcorange!25] (end) at (0,-3.8) {\texttt{print("Go!")}};
  \draw[->,thick] (start) -- (check);
  \draw[->,thick] (check) -- node[above]{yes (True)} (body);
  \draw[->,thick] (body.north) -- ++(0,0.8) -| node[pos=0.25,above]{back to the check} (check.north east);
  \draw[->,thick] (check) -- node[right]{no (False): leave the loop} (end);
\end{tikzpicture}
\end{center}
This prints \texttt{3}, \texttt{2}, \texttt{1}, \texttt{Go!}. The line \texttt{n = n - 1} is the \textbf{change}: without it, \texttt{n} stays 3, the
check is always \texttt{True}, and the loop never ends. That is an \textbf{infinite loop}.

\begin{tip}[Stuck in a loop?]
Press the red \menu{Stop} button in Thonny's toolbar. Then ask two questions:
\textbf{What is checked?} and \textbf{What changes inside the loop, so the check can become False?}
\end{tip}

\subsection{while True and break}
Sometimes it is easier to loop \enquote{forever} and jump out when you are done. \texttt{break} leaves the loop right away.
Inside a function, \texttt{return} leaves the loop \emph{and} the function.
\begin{pycode}[Keep asking until the answer is good]
while True:
    answer = input("Type yes or no: ").lower()   # .lower(): YES becomes yes
    if answer == "yes" or answer == "no":
        break                       # good answer: leave the loop
    print("Please type yes or no.")
print("Thanks!")
\end{pycode}
Checking the player's answer like this is called \textbf{input validation}. Games do it all the time, so a typo never crashes them.
\texttt{answer.isdigit()} is \texttt{True} only if the answer is made of the digits 0--9: then \texttt{int(answer)} is safe.

\subsection{A loop inside a loop: pixel art}
A picture can be a list of strings: one string for each row, one letter for each pixel.
The \textbf{outer loop} goes down the rows. The \textbf{inner loop} goes along one row, pixel by pixel.

\begin{center}
%% alt: A tiny picture 3 columns wide and 2 rows high, made from rows of text. Row 0 is "K.K": black, empty, black. Row 1 is ".O.": empty, orange, empty. K means black, O means orange and a dot means empty. The inner loop moves to the next pixel along a row; the outer loop moves to the next row.
\begin{tikzpicture}[font=\small,>={Stealth[length=3mm]}]
  \fill[black] (0,0) rectangle ++(0.8,-0.8);
  \fill[black] (1.6,0) rectangle ++(0.8,-0.8);
  \fill[kcorange] (0.8,-0.8) rectangle ++(0.8,-0.8);
  \draw (0,0) grid[step=0.8] (2.4,-1.6);
  \node[right] at (2.6,-0.4) {\texttt{"K.K"}\quad row 0};
  \node[right] at (2.6,-1.2) {\texttt{".O."}\quad row 1};
  \draw[->,thick,kcblue] (-0.3,-0.2) -- (-0.3,-1.5) node[midway,left,align=right]{outer loop:\\next row};
  \draw[->,thick,kcvermillion] (0,0.3) -- (2.4,0.3) node[midway,above]{inner loop: next pixel};
  \node[right,align=left] at (6.3,-0.8) {K = black\\O = orange\\. = empty};
\end{tikzpicture}
\end{center}
\begin{pycode}
for row in grid:            # outer loop: one row at a time
    for pixel in row:       # inner loop: one letter at a time
        print(pixel)
\end{pycode}

\begin{exercise}[Guess My Number]
Folder: \texttt{exercises/week04/ex-guess-number}. Open \texttt{starter.py} and press \menu{Run}.
\begin{enumerate}[leftmargin=1.4em]
  \item Guess until you are right. Does the game stop? If not, press \menu{Stop}. Why does it keep going?
  \answerlines{1}
  \item TODO 1: add 1 to \texttt{tries} each time through. TODO 3: \texttt{break} when the result is \texttt{"Correct"}.
  \item Now type \texttt{banana}. Crash! TODO 2: make \texttt{ask\_number()} keep asking until the answer is a whole number from \texttt{low} to \texttt{high}.
  \item Be the tester: try \texttt{banana}, \texttt{99}, \texttt{0} and \texttt{-3}. Does your game survive them all? \tickbox
  \item Check your work with \texttt{check\_my\_work.py}.
\end{enumerate}
\end{exercise}

\begin{exercise}[Pixel Art]
Folder: \texttt{exercises/week04/ex-pixel-art}. Only the top row of the space invader appears.
\begin{enumerate}[leftmargin=1.4em]
  \item TODO 1: wrap the drawing in an outer loop \texttt{for row in grid:}. Move down after each row: \texttt{y = y - size}.
  \item TODO 2: finish \texttt{count\_pixels()} with a loop inside a loop.
  \item Design your own picture (8 rows of 11 letters fit nicely). Plan it here first:
\end{enumerate}
\begin{center}
%% alt: A blank grid of small squares, used to plan a pixel picture.
\begin{tikzpicture}
  \draw[kcgray!70,step=0.6] (0,0) grid (6.6,4.8);
\end{tikzpicture}
\end{center}
Check your work with \texttt{check\_my\_work.py}.
\end{exercise}

\begin{vocab}
  \vocabword{while loop}{A loop that goes around again and again while its condition is True.}
  \vocabword{condition}{The check at the top of a while loop, like \texttt{n > 0}. It is True or False.}
  \vocabword{infinite loop}{A loop that never stops, because nothing changes the condition. Press Stop!}
  \vocabword{break}{Jumps out of a loop right away.}
  \vocabword{input validation}{Checking the player's answer and asking again if it is not OK.}
  \vocabword{nested loop}{A loop inside another loop, like rows and pixels.}
  \vocabword{game loop}{The loop at the heart of a game: get input, update, show, repeat.}
\end{vocab}

\begin{reflection}
When would you use a \texttt{while} loop instead of a \texttt{for} loop?
\answerlines{2}
Name one way a game could be unfair to players. How would you make it fair?
\answerlines{2}
\end{reflection}
